\scanpage{110}
\mspage{55}
% archive sheet number: 109
laquelle induira automatiquement

\[
 \varphi_0p_0\mathbin{\unsure{-|\to}}\mathrm{id},
\]

puisque $\underline{\mathrm{Fibsc}}_0\,S$ est sous-catégorie pleine de $\underline{\mathrm{Fibsc}}\,S$. En fait, on va définir $\varphi p\mathbin{\unsure{-|\to}}\mathrm{id}$, donc

\[
 \boxed{\unsure{\sigma|r|\gamma}_X:\Phi(X)\longrightarrow X}
\]

quand $X$ est fibrée scindée sur $X$ -- une flèche \uline{cartésienne}, donc (\unsure{et on veut|si on veut}) compatible avec les scindages de $\Phi(X)$, et de $X$. En fait, on aura une telle \deleted{\unsure{fonct.}} flèche, fonctorielle en $X$, pour $X$ variable dans $\underline{\mathrm{Fib}}\,S$ (donc sans se borner à des $S$-morphismes cartésiens entre cat. fibrées sur $S$), mais la flèche est alors bien sûr \uline{cartésienne}. (Ça a un sens, sans choisir de scindage sur $X$, et sans qu'il en existe forcément) La définition de cette flèche $\sigma$ fait appel \deleted{\unsure{essentiellement}} \inserted{de façon essentielle} à l'hyp. que $X$ soit \inserted{cat.} \unsure{fibrant|fibrante} sur $S$. Bien sûr, comparée à la flèche \inserted{d'inclusion} évidente

\[
 X\xrightarrow{i_X}\Phi(X)
\]

on va avoir

\[
 \sigma_Xi_X=\mathrm{id}_X.
\]

\margin{[NB. Comme \inserted{on sait que} $i_X\in W$, il en résultera bien $\sigma_X\in W$ comme on le veut)}

On définit $\sigma$ \unsure{fibre|fibres} par \unsure{fibre|fibres},

\[
 \sigma_{X,s}:\Phi(X)_s\overset{\mathrm{def}}{=}{}_{s\backslash}X\longrightarrow X_s,
\]

donc par une \uline{rétraction} de $ {}_{s\backslash}X$ sur la sous-catégorie \inserted{pleine} $X_s$. Un objet de $ {}_{s\backslash}X$ est un couple $(x,u)$, $x\in\mathrm{Ob}\,X$, $u:\deleted{\unsure{p_X(x)}}s\longrightarrow p_X(x)$ ($p_X:X\longrightarrow S$ le morphisme structural de $X$). On pose

\[
 \sigma(x,u)=u^*_X(x)\in\mathrm{Ob}\,X_s,
\]

